% ===== PRÉAMBULE CANONIQUE — à coller VERBATIM en tête de CHAQUE .tex =====
\documentclass[11pt,a4paper]{article}
\usepackage[utf8]{inputenc}\usepackage[T1]{fontenc}\usepackage{lmodern}
\usepackage[french]{babel}
\usepackage{amsmath,amssymb,mathtools}\usepackage{icomma}
\usepackage[version=4]{mhchem}          % \ce{...} : équations & formules
\usepackage{siunitx}\sisetup{locale=FR} % \qty{}{}, \si{}, \ang{}
\usepackage{chemfig}                    % \chemfig{...} : structures développées
\usepackage{xcolor}\usepackage{colortbl}
\usepackage{tikz}\usepackage{pgfplots}\pgfplotsset{compat=1.18}
\usepackage[most]{tcolorbox}\usepackage{enumitem}\usepackage{array,booktabs}
\usepackage[a4paper,margin=2.2cm]{geometry}
\usetikzlibrary{arrows.meta,calc,angles,quotes,positioning,patterns,backgrounds,shapes.geometric}
\definecolor{primary}{RGB}{222,49,99}\definecolor{primarydark}{RGB}{150,22,55}
\definecolor{defcol}{RGB}{0,114,178}\definecolor{methcol}{RGB}{52,92,148}
\definecolor{excol}{RGB}{0,135,90}\definecolor{attncol}{RGB}{200,40,55}
\tcbset{boxbase/.style={enhanced,breakable,boxrule=0.9pt,arc=2.5pt,
  left=3mm,right=3mm,top=2.3mm,bottom=2.3mm,fonttitle=\bfseries,coltitle=white}}
% --- boîtes TUTORIEL (noms EXACTS, ne pas renommer) ---
\newtcolorbox{cle}[1][]{boxbase,colback=defcol!6,colframe=defcol,
  title={\textsc{À retenir}\ifx\relax#1\relax\relax\else\textmd{ : \itshape #1}\fi}}
\newtcolorbox{methode}[1][]{boxbase,colback=methcol!7,colframe=methcol,sharp corners,
  title={\textsc{Méthode}\ifx\relax#1\relax\relax\else\textmd{ --- \itshape #1}\fi}}
\newtcolorbox{exemplebox}[1][]{boxbase,colback=excol!5,colframe=excol,
  title={\textsc{Exemple}\ifx\relax#1\relax\relax\else\textmd{ : \itshape #1}\fi}}
\newtcolorbox{piege}[1][]{boxbase,colback=attncol!6,colframe=attncol,
  title={\textsc{Piège}\ifx\relax#1\relax\relax\else\textmd{ : \itshape #1}\fi}}
\newtcolorbox{remarque}[1][]{boxbase,colback=black!4,colframe=black!45,
  title={\textsc{Remarque}\ifx\relax#1\relax\relax\else\textmd{ : \itshape #1}\fi}}
% --- boîtes EXERCICE / CORRIGÉ (noms EXACTS, auto-comptées) ---
\newtcolorbox[auto counter]{exo}[1][]{boxbase,colback=primary!4,colframe=primary,
  title={\textsc{Exercice}~\thetcbcounter\ifx\relax#1\relax\relax\else\textmd{ --- \itshape #1}\fi}}
\newtcolorbox[auto counter]{corrige}[1][]{boxbase,colback=excol!4,colframe=excol,
  title={\textsc{Corrigé}~\thetcbcounter\ifx\relax#1\relax\relax\else\textmd{ --- \itshape #1}\fi}}
\newcommand{\R}{\mathbb{R}}\newcommand{\N}{\mathbb{N}}\newcommand{\Z}{\mathbb{Z}}\newcommand{\ds}{\displaystyle}
% (un \usepackage{fancyhdr}+en-tête est possible ; il est ignoré côté web)
% =========================================================================
\begin{document}

\begin{exo}[masse, moles et volume aux CNTP]
On donne $M(\ce{O2})=\qty{32}{\gram\per\mole}$, $M(\ce{He})=\qty{4}{\gram\per\mole}$,
$M(\ce{CO2})=\qty{44}{\gram\per\mole}$ ; $V_{\mathrm{m}}=\qty{22.4}{\litre\per\mole}$.
\begin{enumerate}[label=\alph*)]
  \item Quelle masse de \ce{O2} occupe $\qty{11.2}{\litre}$ aux CNTP ?
  \item Quel volume occupent $\qty{8.0}{\gram}$ de \ce{He} aux CNTP ?
  \item Quel volume occupent $\qty{22}{\gram}$ de \ce{CO2} aux CNTP ?
\end{enumerate}
\end{exo}

\begin{exo}[équation d'état hors CNTP]
On prend $R=\num{0,082}$ (unités \si{atm}, \si{\litre}, \si{\kelvin}). Convertis d'abord la
température en kelvins.
\begin{enumerate}[label=\alph*)]
  \item $n=\qty{0.50}{\mole}$ dans $V=\qty{5.0}{\litre}$ à \qty{27}{\degreeCelsius} : calcule $p$.
  \item $n=\qty{2.0}{\mole}$ sous $p=\qty{1.5}{atm}$ à \qty{27}{\degreeCelsius} : calcule $V$.
  \item $p=\qty{2.0}{atm}$, $V=\qty{3.0}{\litre}$, $n=\qty{0.20}{\mole}$ : calcule $T$.
\end{enumerate}
\end{exo}

\begin{exo}[du volume au nombre de molécules]
On rappelle le nombre d'Avogadro $N_A=\qty{6.022e23}{}$ entités par mole ; $V_{\mathrm{m}}=\qty{22.4}{\litre\per\mole}$
aux CNTP.
\begin{enumerate}[label=\alph*)]
  \item Combien de molécules contient $\qty{2.24}{\litre}$ de gaz aux CNTP ?
  \item Combien de molécules contient $\qty{1.0}{\mole}$ de gaz ?
\end{enumerate}
\end{exo}

\begin{exo}[travailler en unités SI ($R=\num{8,314}$)]
On utilise $R=\qty{8.314}{\joule\per\mole\per\kelvin}$, qui impose $p$ en \si{\pascal}, $V$ en
\si{\meter\cubed}, $T$ en \si{\kelvin}.
\begin{enumerate}[label=\alph*)]
  \item Calcule la pression exercée par $n=\qty{1.0}{\mole}$ de gaz dans
        $V=\qty{22.4e-3}{\meter\cubed}$ à $T=\qty{273}{\kelvin}$.
  \item Vérifie que le résultat correspond bien à environ $\qty{1}{atm}$.
\end{enumerate}
\end{exo}

\end{document}
